\section{Thiết kế API}

API của hệ thống được thiết kế theo REST, dùng JSON, đặt dưới tiền tố \texttt{/v1} và dùng quy ước \texttt{camelCase}. Backend NestJS sinh Swagger/OpenAPI từ controller và DTO; khi chạy, Swagger UI có tại \texttt{/docs} và đặc tả JSON có tại \texttt{/docs/json}. Thiết kế API tách \textit{control plane} cho quản trị cấu hình, \textit{data plane} cho runtime evaluation và \textit{monitoring} cho audit/statistics; các endpoint dùng khóa nghiệp vụ ổn định như \texttt{projectKey}, \texttt{flagKey}, \texttt{groupKey} và request ID \texttt{X-Request-Id} để hỗ trợ truy vết.

\subsection{Control plane API}

Control plane quản lý project, flag, rule, group và sample user. Nhóm endpoint này yêu cầu bearer token demo \texttt{demoBearer}; backend resolve actor theo token, kiểm tra quyền theo vai trò \texttt{ADMIN}/\texttt{DEVELOPER}/\texttt{VIEWER}, chuẩn hóa pagination bằng \texttt{items}/\texttt{page}, và ghi audit log cho các mutation quan trọng. Environment chưa có CRUD riêng; dữ liệu environment được dùng qua seed/model và được tham chiếu bằng \texttt{environmentKey} trong evaluation, group config và statistics.

\subsubsection{Project API}

\begin{table}[H]
    \centering
    \scriptsize
    \caption{Đặc tả Project API}
    \label{tab:api-projects}
    \begin{tabular}{@{}p{0.12\textwidth} p{0.34\textwidth} p{0.19\textwidth} p{0.26\textwidth}@{}}
        \toprule
        \textbf{Method} & \textbf{Endpoint} & \textbf{Quyền} & \textbf{Mục đích và dữ liệu chính} \\
        \midrule
        GET & \path|/v1/projects| & \texttt{CONTROL\_PLANE\_READ} & Liệt kê project, trả \texttt{items} và \texttt{page}. Query kế thừa pagination. \\
        POST & \path|/v1/projects| & \texttt{PROJECT\_MANAGE} & Tạo project với \texttt{key}, \texttt{name}, \texttt{description}. Trả \texttt{201 Created}. \\
        GET & \path|/v1/projects/{projectKey}| & \texttt{CONTROL\_PLANE\_READ} & Xem chi tiết project theo khóa nghiệp vụ. \\
        PATCH & \path|/v1/projects/{projectKey}| & \texttt{PROJECT\_MANAGE} & Cập nhật \texttt{name} hoặc \texttt{description}; \texttt{key} là bất biến. \\
        DELETE & \path|/v1/projects/{projectKey}| & \texttt{PROJECT\_MANAGE} & Xóa project ở mức quản trị, trả \texttt{204 No Content}. \\
        \bottomrule
    \end{tabular}
\end{table}

\subsubsection{Feature Flag API}

\begin{table}[H]
    \centering
    \scriptsize
    \caption{Đặc tả Feature Flag API}
    \label{tab:api-feature-flags}
    \begin{tabular}{@{}p{0.12\textwidth} p{0.38\textwidth} p{0.17\textwidth} p{0.24\textwidth}@{}}
        \toprule
        \textbf{Method} & \textbf{Endpoint} & \textbf{Quyền} & \textbf{Mục đích và dữ liệu chính} \\
        \midrule
        GET & \path|/v1/projects/{projectKey}/flags| & \texttt{CONTROL\_PLANE\_READ} & Liệt kê flag chưa soft-delete. Query hỗ trợ \texttt{search}, \texttt{status}, \texttt{lifecycleStatus} và pagination. \\
        GET & \path|/v1/projects/{projectKey}/flags/deleted| & \texttt{CONTROL\_PLANE\_READ} & Liệt kê flag đã soft-delete để có thể khôi phục khi cần. \\
        POST & \path|/v1/projects/{projectKey}/flags| & \texttt{FLAG\_MANAGE} & Tạo flag với \texttt{key}, \texttt{name}, \texttt{description}; backend tạo config mặc định cho environment mặc định. \\
        GET & \path|/v1/projects/{projectKey}/flags/{flagKey}| & \texttt{CONTROL\_PLANE\_READ} & Xem chi tiết một flag, gồm group summary nếu flag đang được gán group. \\
        PATCH & \path|/v1/projects/{projectKey}/flags/{flagKey}| & \texttt{FLAG\_MANAGE} & Cập nhật metadata và config runtime: \texttt{name}, \texttt{description}, \texttt{status}, \texttt{servingMode}, \texttt{killSwitch}. \\
        POST & \path|/v1/projects/{projectKey}/flags/{flagKey}/archive| & \texttt{FLAG\_LIFECYCLE\_MANAGE} & Archive flag; data plane sẽ trả tắt với reason \texttt{FLAG\_ARCHIVED}. \\
        POST & \path|/v1/projects/{projectKey}/flags/{flagKey}/restore| & \texttt{FLAG\_LIFECYCLE\_MANAGE} & Khôi phục flag đã archive về vòng đời active. \\
        DELETE & \path|/v1/projects/{projectKey}/flags/{flagKey}| & \texttt{FLAG\_LIFECYCLE\_MANAGE} & Soft-delete flag, trả \texttt{204 No Content}. \\
        POST & \path|/v1/projects/{projectKey}/flags/{flagKey}/restore-deleted| & \texttt{FLAG\_LIFECYCLE\_MANAGE} & Khôi phục flag đã soft-delete. \\
        \bottomrule
    \end{tabular}
\end{table}

\subsubsection{Rule API}

\begin{table}[H]
    \centering
    \scriptsize
    \caption{Đặc tả Rule API}
    \label{tab:api-rules}
    \begin{tabular}{@{}p{0.12\textwidth} p{0.38\textwidth} p{0.18\textwidth} p{0.23\textwidth}@{}}
        \toprule
        \textbf{Method} & \textbf{Endpoint} & \textbf{Quyền} & \textbf{Mục đích và dữ liệu chính} \\
        \midrule
        GET & \path|/v1/projects/{projectKey}/flags/{flagKey}/rules| & \texttt{CONTROL\_PLANE\_READ} & Liệt kê rule của flag. Query hỗ trợ \texttt{type} và pagination. \\
        PUT & \path|/v1/projects/{projectKey}/flags/{flagKey}/rules| & \texttt{RULE\_MANAGE} & Thay thế toàn bộ rule set bằng mảng \texttt{rules}. Mỗi rule gồm \texttt{type}, \texttt{priority}, \texttt{enabled}, \texttt{parameters}. \\
        \bottomrule
    \end{tabular}
\end{table}

\subsubsection{Flag Group API}

\begin{table}[H]
    \centering
    \scriptsize
    \caption{Đặc tả Flag Group API}
    \label{tab:api-flag-groups}
    \begin{tabular}{@{}p{0.12\textwidth} p{0.38\textwidth} p{0.18\textwidth} p{0.23\textwidth}@{}}
        \toprule
        \textbf{Method} & \textbf{Endpoint} & \textbf{Quyền} & \textbf{Mục đích và dữ liệu chính} \\
        \midrule
        GET & \path|/v1/projects/{projectKey}/groups| & \texttt{CONTROL\_PLANE\_READ} & Liệt kê group. Query hỗ trợ \texttt{search}, \texttt{environmentKey} và pagination. \\
        POST & \path|/v1/projects/{projectKey}/groups| & \texttt{GROUP\_MANAGE} & Tạo group với \texttt{key} và \texttt{name}. \texttt{key} là bất biến trong project. \\
        PATCH & \path|/v1/projects/{projectKey}/groups/{groupKey}| & \texttt{GROUP\_MANAGE} & Đổi \texttt{name} của group. \\
        DELETE & \path|/v1/projects/{projectKey}/groups/{groupKey}| & \texttt{GROUP\_MANAGE} & Xóa group, trả \texttt{204 No Content}. \\
        PUT & \path|/v1/projects/{projectKey}/groups/{groupKey}/config| & \texttt{GROUP\_KILL\_SWITCH} & Cập nhật \texttt{killSwitch} cho một \texttt{environmentKey}. \\
        PUT & \path|/v1/projects/{projectKey}/flags/{flagKey}/group| & \texttt{GROUP\_ASSIGN} & Gán flag vào group bằng body \texttt{groupKey}; trả lại flag response mới. \\
        DELETE & \path|/v1/projects/{projectKey}/flags/{flagKey}/group| & \texttt{GROUP\_ASSIGN} & Gỡ flag khỏi group; trả lại flag response mới. \\
        \bottomrule
    \end{tabular}
\end{table}

\subsubsection{Sample User API}

\begin{table}[H]
    \centering
    \scriptsize
    \caption{Đặc tả Sample User API}
    \label{tab:api-sample-users}
    \begin{tabular}{@{}p{0.12\textwidth} p{0.36\textwidth} p{0.20\textwidth} p{0.23\textwidth}@{}}
        \toprule
        \textbf{Method} & \textbf{Endpoint} & \textbf{Quyền} & \textbf{Mục đích và dữ liệu chính} \\
        \midrule
        GET & \path|/v1/projects/{projectKey}/sample-users| & \texttt{CONTROL\_PLANE\_READ} & Liệt kê sample user. Query hỗ trợ \texttt{search}, \texttt{role} và pagination. \\
        POST & \path|/v1/projects/{projectKey}/sample-users| & \texttt{SAMPLE\_USER\_MANAGE} & Tạo sample user với \texttt{displayName}, \texttt{targetingKey}, \texttt{userId}, \texttt{roles}, \texttt{attributes}. \\
        DELETE & \path|/v1/projects/{projectKey}/sample-users/{targetingKey}| & \texttt{SAMPLE\_USER\_MANAGE} & Xóa sample user theo \texttt{targetingKey}, trả \texttt{204 No Content}. \\
        \bottomrule
    \end{tabular}
\end{table}

\subsection{Data plane API}

Data plane chỉ có một endpoint chính: \texttt{POST /v1/evaluate}. Endpoint này được đánh dấu public trong controller, không dùng bearer token quản trị, và được tối ưu cho tính xác định, fail-closed và an toàn khi thiếu cấu hình. Đây là API mà Demo Application gọi thông qua \texttt{@ffp/js-sdk} để quyết định có hiển thị một tính năng hay không.

\begin{table}[H]
    \centering
    \scriptsize
    \caption{Đặc tả Evaluation API}
    \label{tab:api-evaluation}
    \begin{tabular}{@{}p{0.12\textwidth} p{0.24\textwidth} p{0.20\textwidth} p{0.35\textwidth}@{}}
        \toprule
        \textbf{Method} & \textbf{Endpoint} & \textbf{Quyền} & \textbf{Mục đích và dữ liệu chính} \\
        \midrule
        POST & \path|/v1/evaluate| & Public & Trả quyết định runtime cho một \texttt{projectKey}, \texttt{flagKey}, \texttt{environmentKey} tùy chọn và \texttt{context}. Response luôn có dạng evaluation result với \texttt{enabled}, \texttt{variant}, \texttt{reason}. \\
        \bottomrule
    \end{tabular}
\end{table}

Request body của Evaluation API gồm khóa cấu hình và context runtime:

\begin{lstlisting}[
    caption={Ví dụ request Evaluation API},
    label={lst:evaluation-api-request}
]
{
  "projectKey": "demo-project",
  "environmentKey": "production",
  "flagKey": "new-checkout",
  "context": {
    "targetingKey": "demo-user-beta",
    "userId": "demo-user-beta",
    "roles": ["beta-tester"],
    "attributes": {
      "plan": "pro",
      "country": "VN"
    }
  }
}
\end{lstlisting}

Response body trả lại quyết định đã được giải thích bằng reason code:

\begin{lstlisting}[
    caption={Ví dụ response Evaluation API},
    label={lst:evaluation-api-response}
]
{
  "projectKey": "demo-project",
  "flagKey": "new-checkout",
  "enabled": true,
  "variant": "on",
  "reason": "ROLE_MATCH",
  "matchedRuleId": "rule_123"
}
\end{lstlisting}

Nếu project, environment, flag hoặc config theo environment không tồn tại, API không trả lỗi \texttt{404} kiểu control plane mà trả kết quả an toàn: \texttt{enabled = false}, \texttt{variant = "off"}, \texttt{reason = "NOT\_FOUND"}. Nếu request hợp lệ nhưng thiếu context cần thiết cho percentage rollout, engine trả \texttt{INVALID\_CONTEXT}. Nếu có lỗi bất ngờ trong đường evaluate, service bắt lỗi và trả \texttt{ERROR} theo nguyên tắc fail-closed. Các reason code hiện có gồm \texttt{GLOBAL\_ON}, \texttt{FLAG\_DISABLED}, \texttt{FLAG\_ARCHIVED}, \texttt{GROUP\_KILL\_SWITCH}, \texttt{KILL\_SWITCH}, \texttt{USER\_ALLOWLIST}, \texttt{ROLE\_MATCH}, \texttt{PERCENTAGE\_ROLLOUT}, \texttt{DEFAULT\_OFF}, \texttt{NOT\_FOUND}, \texttt{INVALID\_CONTEXT} và \texttt{ERROR}.

Thứ tự ưu tiên của data plane là: archived flag, disabled config, group kill switch, flag kill switch, global on, ordered enabled rules và default off. Vì vậy một rule rollout không thể vô tình bật flag khi flag đã bị archive, disabled hoặc bị kill switch.

\subsection{Monitoring API}

Monitoring API gồm audit/history và evaluation statistics. Audit log là nguồn append-only để truy vết actor, action, target, request ID và snapshot trước/sau; flag history là projection từ audit log, không tạo nguồn sự thật thứ hai. Statistics chỉ lưu dữ liệu aggregate theo flag, environment, reason, enabled/disabled và bucket giờ UTC, không lưu raw evaluation context hay định danh người dùng cuối. Ngoài ra backend có health probe công khai \path|GET /v1/health|.

\subsubsection{Audit Log API}

\begin{table}[H]
    \centering
    \scriptsize
    \caption{Đặc tả Audit Log và Flag History API}
    \label{tab:api-audit-history}
    \begin{tabular}{@{}p{0.12\textwidth} p{0.38\textwidth} p{0.19\textwidth} p{0.22\textwidth}@{}}
        \toprule
        \textbf{Method} & \textbf{Endpoint} & \textbf{Quyền} & \textbf{Mục đích và dữ liệu chính} \\
        \midrule
        GET & \path|/v1/projects/{projectKey}/audit-logs| & \texttt{CONTROL\_PLANE\_READ} & Liệt kê audit log của project. Query hỗ trợ \texttt{targetType}, \texttt{targetKey}, \texttt{actor}, \texttt{action}, \texttt{from}, \texttt{to} và pagination. \\
        GET & \path|/v1/projects/{projectKey}/flags/{flagKey}/history| & \texttt{CONTROL\_PLANE\_READ} & Lịch sử cấu hình của một flag, được projection từ audit log liên quan đến flag, config và rule. Query hỗ trợ pagination và \texttt{sort=createdAt}. \\
        \bottomrule
    \end{tabular}
\end{table}

\subsubsection{Evaluation Statistics API}

\begin{table}[H]
    \centering
    \scriptsize
    \caption{Đặc tả Evaluation Statistics API}
    \label{tab:api-statistics}
    \begin{tabular}{@{}p{0.12\textwidth} p{0.36\textwidth} p{0.19\textwidth} p{0.24\textwidth}@{}}
        \toprule
        \textbf{Method} & \textbf{Endpoint} & \textbf{Quyền} & \textbf{Mục đích và dữ liệu chính} \\
        \midrule
        GET & \path|/v1/projects/{projectKey}/stats/flags| & \texttt{CONTROL\_PLANE\_READ} & Thống kê tổng hợp theo các flag trong project. Query hỗ trợ \texttt{environmentKey}, \texttt{from}, \texttt{to}, pagination và sort theo \texttt{totalEvaluations}, \texttt{enabledCount}, \texttt{disabledCount}, \texttt{flagKey}. \\
        GET & \path|/v1/projects/{projectKey}/flags/{flagKey}/stats| & \texttt{CONTROL\_PLANE\_READ} & Thống kê chi tiết một flag: tổng lượt evaluate, số enabled/disabled, phần trăm enabled, breakdown theo reason và bucket theo giờ. Query hỗ trợ \texttt{environmentKey}, \texttt{from}, \texttt{to}. \\
        \bottomrule
    \end{tabular}
\end{table}

\subsection{Validation và mô hình lỗi}

Validation được thực hiện ở biên API bằng DTO, \texttt{class-validator}, \texttt{class-transformer} và global \texttt{ValidationPipe}. Pipe bật \texttt{whitelist}, \texttt{forbidNonWhitelisted} và \texttt{transform}, vì vậy field ngoài DTO bị từ chối, query string như \texttt{limit} được chuyển sang số, và lỗi validation được chuẩn hóa trước khi vào service. Các khóa nghiệp vụ như \texttt{projectKey}, \texttt{flagKey}, \texttt{groupKey} và \texttt{environmentKey} dùng regex:

\begin{lstlisting}[
    caption={Định dạng key nghiệp vụ},
    label={lst:key-regex}
]
^[a-z0-9][a-z0-9-]{1,62}[a-z0-9]$
\end{lstlisting}

Quy tắc này tương ứng độ dài 3--64 ký tự, chỉ gồm chữ thường, chữ số và dấu gạch ngang, không bắt đầu hoặc kết thúc bằng dấu gạch ngang. Các DTO còn giới hạn độ dài \texttt{name} tối đa 120 ký tự, \texttt{description} tối đa 1000 ký tự, pagination \texttt{limit} từ 1 đến 100, \texttt{offset} không âm, thời gian \texttt{from}/\texttt{to} theo ISO-8601 và enum phải thuộc giá trị hợp lệ.

Control-plane error response có dạng thống nhất:

\begin{lstlisting}[
    caption={Mẫu error response của control plane},
    label={lst:api-error-response}
]
{
  "code": "VALIDATION_ERROR",
  "message": "Request validation failed.",
  "details": [
    {
      "field": "projectKey",
      "message": "projectKey must be a valid 3-64 character key."
    }
  ],
  "requestId": "req_123"
}
\end{lstlisting}

Các mã lỗi chuẩn gồm \texttt{VALIDATION\_ERROR}, \texttt{UNAUTHORIZED}, \texttt{FORBIDDEN}, \texttt{NOT\_FOUND}, \texttt{CONFLICT} và \texttt{INTERNAL\_ERROR}. \texttt{UNAUTHORIZED} dùng khi thiếu hoặc sai bearer token demo; \texttt{FORBIDDEN} dùng khi role đã xác thực nhưng không có permission; \texttt{CONFLICT} dùng cho xung đột như khóa đã tồn tại; \texttt{INTERNAL\_ERROR} che giấu chi tiết lỗi nội bộ nhưng vẫn log kèm request ID. Riêng Evaluation API có thêm mô hình lỗi domain bằng reason code trong response 200 để giữ fail-closed, ví dụ \texttt{NOT\_FOUND}, \texttt{INVALID\_CONTEXT} hoặc \texttt{ERROR} đều trả \texttt{enabled = false} và \texttt{variant = off} thay vì để client tự suy diễn.
